\documentclass[11pt]{article}

% Language setting
% Replace `english' with e.g. `spanish' to change the document language
\usepackage[english]{babel}

% Set page size and margins
% Replace `letterpaper' with `a4paper' for UK/EU standard size
\usepackage[letterpaper,top=2cm,bottom=2cm,left=3cm,right=3cm,marginparwidth=1.75cm]{geometry}

% Useful packages
\usepackage{amsmath}
\usepackage{graphicx}
\usepackage[colorlinks=true, allcolors=blue]{hyperref}
\usepackage{float}
\usepackage{natbib}

\title{Your Paper}
\author{You}

\begin{document}
\maketitle

\section{Introduction}

Cloud services are very popular nowadays to replace local servers, which can significantly reduce the cost of maintaining the servers as well as human resources. However, if there are security misconfigurations, it will put enterprises in danger of incidents such as data leaks. Although cloud platforms have provided detailed secure configuration guidelines to help engineers use their services, the guidelines contain a vast amount of complex information, causing low search efficiency when engineers handle these large documents. This project aims to implement Naive RAG and Hybrid RAG to retrieve context from the documents, and uses an open-source, instruction-tuned LLM to generate actionable and precise security guidelines for practitioners. Finally, it evaluates the performance of both RAG variants using an LLM-as-a-Judge framework.

\section{Literature Review \& System Positioning}

\subsection{Analysis of Existing Cybersecurity RAG Frameworks}

The paper on CyberBOT demonstrates ways to solve the reliability issue in cybersecurity education by using ontology-grounded RAG. It provides an inspiration that a standard RAG system can easily make mistakes in areas that require high accuracy, such as cybersecurity. Therefore, a strict structured control over the retrieved content is necessary. Besides,the CyberRAG study has also shown that using agentic RAG can automatically classify cyber-attacks and generate reports. An advanced RAG system should not only be able to read context, but also match the most accurate content from the knowledge base according to practitioner needs.

\subsection{Methodological Inspiration \& Positioning}

This project focuses on the secure configuration for the Azure cloud rather than conceptual cybersecurity knowledge. The Azure cloud is highly popular across enterprises of all sizes, regardless of whether they are large, medium, or small. However, because its official documentation is filled with a large number of precise proper nouns and complex configuration rules, such as Microsoft Entra ID, Azure Key Vault, TLS 1.2, HTTPS-only models, and NSG rules, practitioners who are unfamiliar with these details may leave the entire system in an insecure state, regardless of how secure the default design might be. Therefore, implementing a RAG system to search for Azure cloud security information from official documents can effectively help practitioners set up their cloud projects securely.

Naive RAG is the traditional way to retrieve data from a knowledge base. It searches for relevant information by calculating semantic similarity. However, Naive RAG can get confused by Azure services that are totally different but have similar names during calculation, such as mistaking Identity Server's HTTPS settings for Azure App Service settings. This project has decided to implement a hybrid RAG structure to see whether it can provide better results. The design of the hybrid RAG uses dense embeddings to extract the semantics of practitioner queries, combined with a BM25 keyword retriever to perform exact token matching for Azure terms. To compare both RAG systems, DeepEval was used to evaluate system performance by calculating the three core metrics, rather than using RAGAS. This study also researches whether the hybrid structure can make a substantial contribution to the quality of generation for the Azure cloud secure configuration topic.

\section{Methodology \& System Design}

\subsection{Knowledge Base Setup \& Data Pipeline}

The Azure cloud provides well-structured documents for guiding users to implement configurations in their cloud projects. This project connected to the Azure open-source repository by using the GitHub API. By scanning the entire directory, all markdown documents containing the keywords "security" and "baseline" were filtered. Eventually, 323 high-quality official security technical archives were downloaded and stored locally for further processing.The downloaded documents contain many heading levels with markdown "\#", "\#\#" and "\#\#\#", which can influence the retrieval results. Therefore, this project used the "MarkdownHeaderTextSplitter" component from LangChain to perform the first layer of structured segmentation, and then used the "RecursiveCharacterTextSplitter" to do the second layer of segmentation. A chunk size of 1000 was used for the segmentation as it can fully include a well-structured Azure document model, with a chunk overlap of 150 to ensure that the end of the previous block is sufficiently related to the beginning of the next block. To reduce noise from the documents, a blacklist filter was also implemented to remove irrelevant sections, such as "next steps", "feedback", and "see also." After the cleaning and filtering, a total number of 6160 high-quality text blocks were reserved and made ready for the next process.

\subsection{Architectural Variations of Retrieval Module}

For the embedding model and vector store selection, this project has chosen "BAAI/bge-small-en-v1.5" and "Chroma DB" as the primary technologies. "BAAI/bge-small-en-v1.5" ranks highly on the MTEB leaderboard and is known as a model that can accurately extract the semantics of a query.It has a lower vector dimension and an extremely fast calculation speed, and it also works very well in an English environment like Azure documents, making these features highly suitable for this project. As for "Chroma DB", it is a vector database specially designed for AI applications, which is open-source, free, and easy to set up without complex configurations. Since the documents are segmented by LangChain's tools, the Chroma DB API can connect with the text segmentation and search functions easily.

As mentioned previously, Naive RAG and Hybrid RAG were both implemented for performance comparison. For the Naive RAG, the system converts the query into a vector when practitioners submit a query, and uses cosine similarity to retrieve the top-3 nearest text chunks from the Chroma DB. This RAG variant highly depends on semantic similarity, representing the most traditional and basic way to implement a RAG system. The documents were also loaded into the "BM25Retriever" from LangChain to build a sparse index based on traditional word frequency statistics. The system then uses "EnsembleRetriever" to integrate the Naive RAG and BM25 RAG via a hybrid retrieval approach, applying an equal weight of 0.5 : 0.5. After the integration, the system selects the top-3 best text chunks to be subsequently fed into the LLM.

\subsection{Generation Module \& Prompt Engineering}

This project selects "Meta-Llama-3-8B-Instruct" as the system generator. This model is instruction-tuned and optimized with RLHF, giving it a strong prompt-following capability to strictly adhere to instructions, specifically, to generate exactly three concise guidelines based only on the provided context, thereby ensuring the precision and refinement of the output. When handling long Azure documents exceeding thousands of words, it will not be distracted by irrelevant text or prone to over-generation. This model was also trained on vast technical documentation, extensive codebases, and cybersecurity knowledge, making it highly suitable for understanding structured cloud security code and domain-specific technical terms.

To position the model as a senior Azure cloud security architect, this project has designed a rigorous system prompt. The prompt included a closed-world assumption that required the model to only answer the question based on the provided top-3 text chunks. If it cannot find the answer from the top-3 text chunks, it must respond with "I do not have enough information" to ensure no hallucination exists. When the model gives the final answer, the response must contain 3 concise and actionable guidelines that practitioners can follow.

\section{Evaluation Setup \& Metric Framework}

\subsection{Evaluation Dataset Setup}

There are 10 practitioner-style queries related to Azure cloud security were deigned for the evaluation purpose. These 10 queries included critical security scenarios in real world cloud-based operations and maintenance. For instance, monitoring abnormal logins with Microsoft Entra ID, key management and rotation in Azure Key Vault, preventing SQL injection and blocking public network access in Azure SQL Repository, and enforcing HTTPS and TLS 1.2+ security configurations for Azure App Service. The evaluation mode is reference-free that does not require any ground-truth, and rating is based entirely on the logical relationship between the context retrieved by RAG and the generated answer.

\subsection{DeepEval}

The evaluation logic of the RAGAS framework highly relies on the judge model's identification of vagueness and abstract reasoning, especially in reference-free mode. When handling domain-specific scenarios containing extremely precise hardware commands, architecture configurations, and Azure cloud deployment parameters, RAGAS becomes overgeneralized due to a lack of criteria-based granularity. official Azure documentation contains many non-semantic elements, such as title syntax, code comments, and NOTE blocks. These contexts will easily cause the ratings to show an inexplicable and drastic decline due to the algorithm design of RAGAS, making it unsuitable for the deep error analysis of this project.

DeepEval features its own academically grounded evaluation paradigm called the "G-Eval Framework". Unlike RAGAS's fixed black-box rating system, G-Eval allows engineers to customize chain-of-thought evaluation steps for a specific domain, such as Azure cloud security configuration guidelines. This enables the judge model to act like a human security expert, assessing the security of the generated answer step by step based on objective configuration logic. When evaluating faithfulness, DeepEval's calculation method automatically breaks down the guidelines generated by the LLM into independent claims and performs contradiction matrix matching with the retrieved context. This method allows the judge model to keenly capture hallucinations. For instance, the LLM may fabricate a seemingly reasonable but actually non-existent Azure parameter. This zero-tolerance capability is an advantage that RAGAS cannot provide in cybersecurity architecture assessments.

\subsection{Three Core Evaluation Metrics}

There are three core evaluation metrics used to evaluate the RAG system. The answer relevancy metric is used to check whether the generated guidelines truly answer the practitioner's query by analyzing the output's conciseness and accuracy. If there is any irrelevant content or redundancies, the score will be reduced. For faithfulness, it is used to check whether the generated guidelines are fully grounded in the context. The judge model strictly identifies any false Azure parameters, commands, or subjective fabrications that were not mentioned in the context, which ensures that the output is technically reliable. Sometimes the RAG system may retrieve irrelevant noise, such as descriptions of unrelated Azure services, Markdown image tags, or note blocks. The contextual relevancy metric calculates the proportion of truly valuable content within the retrieved top-3 chunks, and gives a high score if the proportion is high.

\subsection{Implementation of Custom Judge LLM}

This project utilizes the Llama-3.3-70B-Versatile model as the evaluator model. According to academic norms of the LLM-as-a-Judge paradigm, the parameter scale of the judge model should be significantly larger than that of the generation model. UUtilizing a 70B large model to evaluate an 8B small model ensures that the ratings possess sufficient objectivity and authoritative review capabilities. This substantial parameter advantage also prevents hallucinations when processing the Chain-of-Thought protocols of G-Eval. Most importantly, this model is accessible via providers like Groq without API costs, making it highly budget-friendly and suitable for this project.

\section{Experimental Results \& Deep Error Analysis}

\begin{figure}[H]
\centering
\includegraphics[width=1\linewidth]{figures/overall_performance_comparison.png}
\caption{Overall Performance Comparison}
\label{fig:figures/overall_performance_comparison}
\end{figure}

\begin{figure}[H]
\centering
\includegraphics[width=1\linewidth]{figures/query_performance_trends.png}
\caption{Query Performance Trends}
\label{fig:figures/query_performance_trends}
\end{figure}

\subsection{Quantitative Results Summary}

Figure 1 demonstrates the average scores of the two RAG variants across 10 queries.The score differences in Answer Relevancy and Faithfulness are very close; in fact, Faithfulness achieves nearly the same score, and the Naive RAG score in Answer Relevancy is only 0.058 higher than Hybrid RAG. This indicates that the prompt engineering for the generator model is a complete success, demonstrating that the model is highly obedient to the instruction that answers must be strictly based on the provided context. However, Naive RAG performs better in Contextual Relevancy. It scores nearly 0.09 higher than Hybrid RAG, showing that BM25 brings more noise to the system, which introduces irrelevant context.

Figure 2 shows the scores of both RAG variants across each query for each metric. For Naive RAG, it displays a perfect straight line at 1.0 in the Answer Relevancy metric, showing that all the answers generated by the LLM perfectly addressed the questions. It also performs well in Faithfulness, achieving a nearly stable straight line across most queries. Scores dropped to around 0.8 only in two queries, Q5 and Q7. The judge noted that these drops occurred because the evidence in the context was not direct enough and the output contained minor reasoning discrepancies. This could be caused by the synthesis capabilities of the model. However, there are strong up-and-down fluctuations in Contextual Relevancy; although it reaches a score of 1.0 in one query, it drops to around 0.5 in another. The judge explained that the context contained substantial noise, including irrelevant formatting and information, which made the content less focused. For Hybrid RAG, the performance on the Answer Relevancy metric is slightly worse than that of Naive RAG. Three queries failed to reach a score of 1.0, with one dropping to around 0.6, because the answers were less precise and occasionally off-topic. However, this RAG variant performs exceptionally well in Faithfulness, with only one query failing to score 1.0. The judge commented that the response contained "a minor implication not explicitly supported by the context." Hybrid RAG exhibits similar strong fluctuations in Contextual Relevancy to Naive RAG, as the retrieved context contains a large amount of irrelevant content and excessive background noise.

\subsection{Deep Error Analysis \& Failure Modes}

Naive RAG performs well in Answer Relevancy, but has a slightly dropped score in Faithfulness. The logs from the judge model show that the retrieved configuration evidence from Naive RAG is not direct enough, constraining the large model's synthesis capabilities when trying to merge and integrate these indirect technical clauses, thereby causing minor reasoning discrepancies. This proves that Naive RAG may suffer from reasoning failures if it relies solely on semantic similarity to retrieve context, especially when handling queries that require multiple discontinuous steps or highly accurate reasoning. Because Naive RAG overemphasizes overall semantic similarity, it inadvertently retrieves context cluttered with irrelevant formatting, such as note blocks and image syntax, as well as background noise. These non-core elements influence the model's attention, significantly decreasing its focus. This represents the "search noise interference" limitation of Naive RAG, leading to a lower score in Contextual Relevancy.

The overall performance of Hybrid RAG is slightly worse than that of Naive RAG in this evaluation. This could be due to a lack of semantic understanding of the security architecture, as BM25 only performs literal token matching. Thus, the system retrieved off-topic context that shares identical keywords, leading to reduced accuracy in the security configuration guidelines generated by the LLM and causing them to deviate from the topic. BM25 is very easy to trigger entity semantic drift in scenarios with an extremely high density of proper nouns like Azure cloud, which incorporates irrelevant material into the context and eventually influences the response quality. As highlighted by the judge model, one of the answers from the Hybrid RAG contains "a minor implication not explicitly supported by the context". This could be due to the sparse search of BM25 bringing substantial background noise, including negative examples or irrelevant policies, such as the Soft Delete 14-day retention policy. The model produced false suggestive derivatives when processing this noisy information, causing latent hallucinations. Although Hybrid RAG can recognize proper nouns, it still shares a similar issue with Naive RAG regarding Contextual Relevancy, showing that it injects a larger amount of search noise into the system while providing accurate token matching.

Based on the results, Naive RAG seems to perform better than Hybrid RAG in Azure cloud configuration scenarios. However, different test cases and architectural configurations could significantly influence the outcomes; therefore, adjusting these settings could lead to different results. For example, adjusting the weights in the Ensemble Retriever from 0.5:0.5 to 0.7:0.3 can allow the system to primarily rely on dense vector semantic control, utilizing BM25 only for minor precision enhancement of proper nouns. This adjustment can significantly reduce the non-core background noise introduced by BM25. Although this project implemented a blacklist to clean the document context, there is still substantial noise that severely dilutes the semantics; if there were a better way to clean and filter the data, both RAG variants might achieve better performance.

\section{Trustworthiness, Ethics \& Governance}

\subsection{Trustworthiness \& Bias}

Although this RAG system is very helpful in guiding people to use cloud services, it poses a risk of providing outdated information.Cloud services update very frequently; if the knowledge base is not updated simultaneously, the system may suggest secure configurations that are already disused or contain vulnerabilities. Besides, he official documents always suggest the best configurations to the user without considering their budget. The system may ignore the needs of small and medium-sized enterprises with limited budgets or organizations with scarce resources, recommending high-cost advanced security services, such as the advanced Defender plan, which leads to technology exclusivity.

\subsection{Transparency}

The RAG system successfully retrieves context from the official documents and provides it to the LLM. However, the top-3 guidelines do not mark their reference sources, which may reduce trust among practitioners. Also, it is difficult for engineers to verify how the AI chooses the 3 guidelines, as the internal reasoning logic is hidden from the user when integrating the documents.

\subsection{Governance \& Human-in-the-Loop}

To solve the issues mentioned above, it is recommended to create an automated knowledge update pipeline, allowing the system to frequently extract the latest official documents from GitHub, ensuring timeliness. Besides, disclaimers and risk warnings should also be included in the generated guidelines, reminding practitioners to conduct a secondary evaluation of the guidelines based on their specific IT environment. Last but not least, the commands generated by the LLM must not be authorized to execute automatically; they must always be reviewed by certified security experts before deployment to the production environment.

\section{Conclusion}

This project has successfully implemented Naive RAG and Hybrid RAG to generate useful Azure secure configuration guidelines for practitioners. Both RAG variants were also evaluated using the DeepEval framework and an LLM-as-a-Judge, producing excellent results. It show that the Llama-3 model demonstrates high faithfulness and anti-hallucination capabilities under strict prompt constraints, stably generating reliable security guidelines. Although the BM25 retriever in the Hybrid structure is skilled at extracting proper nouns, it still introduced irrelevant noise and semantic drift, causing its overall score to be unexpectedly lower than Naive RAG. In the highly specialized field of cloud security, simply overlaying retrieval methods cannot guarantee improved performance. The most critical challenge for optimizing RAG systems in the future is accurately filtering out technical noise and identifying entity boundaries.

\bibliographystyle{alpha}
\bibliography{sample}

\end{document}
